% Enam solusi asli untuk kejadian soal yang tidak mempunyai tautan solusi sumber.
% Bukan bagian dari sumber Brenner dan tidak boleh dinisbahkan kepada Brenner/Wikiversity.
% Lisensi: CC BY-SA 4.0.
% Provenans: disusun oleh OpenAI Codex gpt-5.6-sol, Ultra, atas arahan pengguna.

\section{Solusi asli untuk enam kekosongan pada bank ujian}
\label{o011-exam-original-repairs}

Bagian ini memuat tepat enam solusi yang tidak tersedia pada permukaan sumber
yang dibekukan. Setiap solusi diberi ID terpisah dan secara eksplisit ditandai
sebagai materi asli. Teks soal tetap berasal dari formulir peserta resmi;
solusinya tidak dinisbahkan kepada pengarang atau Wikiversity.

\subsection*{Ujian 1, soal tanpa solusi sumber (O011-EXAM01-ORIG-SOL-01)}
\label{o011-exam01-orig-sol-01}

\paragraph{Soal.}
Deskripsikan sebanyak mungkin manifold berbatas berdimensi satu yang
terhubung.

\paragraph{Solusi asli.}
Dengan konvensi manifold Hausdorff dan berbasiskan terhitung yang dipakai
dalam edisi ini, klasifikasi hingga difeomorfisme adalah
\[
 S^1,\qquad \mathbb R,\qquad [0,\infty),\qquad [0,1].
\]
Alasannya dapat dibaca dari banyaknya titik batas dan kekompakan. Jika
batasnya kosong, manifold satu dimensi terhubung adalah lingkaran apabila
kompak dan garis nyata apabila tidak kompak. Jika batasnya mempunyai dua titik,
manifoldnya difeomorfik dengan interval kompak $[0,1]$. Jika batasnya mempunyai
satu titik, manifoldnya difeomorfik dengan setengah garis $[0,\infty)$.
Interval $[0,1)$ tidak memberikan tipe baru karena
$x\mapsto x/(1-x)$ adalah difeomorfisme ke $[0,\infty)$; semua interval terbuka
tak kosong difeomorfik dengan $\mathbb R$. Jadi daftar tersebut juga menjelaskan
mengapa tidak ada tipe tambahan di bawah hipotesis standar.

\subsection*{Ujian 3, soal tanpa solusi sumber (O011-EXAM03-ORIG-SOL-01)}
\label{o011-exam03-orig-sol-01}

\paragraph{Soal.}
Katakan sesuatu yang cerdas tentang pita M\"obius.

\paragraph{Solusi asli.}
Pita M\"obius dapat dibangun sebagai hasil bagi
\[
 \bigl([0,1]\times[-1,1]\bigr)/\bigl((0,t)\sim(1,-t)\bigr).
\]
Ia merupakan manifold mulus kompak berdimensi dua dengan batas; kedua sisi
vertikal direkatkan dengan pembalikan arah. Dua ruas horizontal bergabung
menjadi satu komponen batas. Garis tengah $[0,1]\times\{0\}$ menjadi sebuah
lingkaran dan merupakan retrak deformasi, sehingga pita M\"obius mempunyai
tipe homotopi $S^1$.

Pita ini tidak dapat diorientasikan. Jika suatu basis berorientasi diangkut
sekali mengelilingi garis tengah, identifikasi $(0,t)\sim(1,-t)$ membalik satu
arah basis dan mengembalikannya dengan orientasi berlawanan. Penutup ganda
berorientasinya adalah silinder. Secara setara, pita M\"obius adalah ruang
total bundel garis nyata tak trivial di atas $S^1$; fakta bahwa ia tertanam
sebagai permukaan di $\mathbb R^3$ tidak membuatnya berorientasi.

\subsection*{Ujian 5, soal tanpa solusi sumber (O011-EXAM05-ORIG-SOL-01)}
\label{o011-exam05-orig-sol-01}

\paragraph{Soal.}
Jelaskan peran teorema tentang pemetaan implisit dalam pengembangan teori
manifold.

\paragraph{Solusi asli.}
Teorema pemetaan implisit adalah mesin lokal yang mengubah persamaan menjadi
koordinat. Jika $F\colon U\subseteq\mathbb R^n\to\mathbb R^r$ mulus dan
$D F_p$ surjektif, maka di sekitar $p$ terdapat koordinat mulus yang di
dalamnya $F$ hanyalah proyeksi pada $r$ koordinat. Akibatnya, himpunan aras
regular
\[
 F^{-1}(F(p))
\]
secara lokal adalah $\mathbb R^{n-r}\times\{0\}$ dan karena itu merupakan
submanifold berdimensi $n-r$. Ruang singgungnya diberikan secara langsung oleh
\[
 T_pF^{-1}(F(p))=\ker D F_p.
\]

Hasil ini memasok banyak contoh dasar: hipermuka sebagai aras regular satu
fungsi, grup matriks yang ditentukan oleh persamaan regular, dan grafik solusi
lokal persamaan. Ia juga menjamin bahwa definisi abstrak manifold cocok dengan
gambaran sebagai ruang solusi lokal. Bersama teorema fungsi invers, ia
menyediakan koordinat, membuktikan kemulusan perubahan koordinat yang sesuai,
dan mendasari teorema rank konstan serta kriteria submanifold. Dengan demikian,
teorema itu menghubungkan kalkulus multivariabel dengan teori manifold.

\subsection*{Ujian 7, soal tanpa solusi sumber (O011-EXAM07-ORIG-SOL-01)}
\label{o011-exam07-orig-sol-01}

\paragraph{Soal.}
Hitung luas permukaan bola satuan.

\paragraph{Solusi asli.}
Parametrisasi bola satuan $S^2\subseteq\mathbb R^3$, kecuali garis bujur dan
kedua kutub yang tidak memengaruhi integral, dengan
\[
 \Phi(\theta,\varphi)
 =(\sin\varphi\cos\theta,\sin\varphi\sin\theta,\cos\varphi),
 \quad 0<\theta<2\pi,\quad0<\varphi<\pi.
\]
Vektor turunannya memenuhi
\[
 \left\|\partial_\theta\Phi\times\partial_\varphi\Phi\right\|
 =\sin\varphi.
\]
Maka luasnya adalah
\[
 \operatorname{luas}(S^2)
 =\int_0^{2\pi}\int_0^\pi\sin\varphi\,d\varphi\,d\theta
 =2\pi\,[ -\cos\varphi ]_0^\pi
 =4\pi.
\]

\subsection*{Ujian 9, soal tanpa solusi sumber (O011-EXAM09-ORIG-SOL-01)}
\label{o011-exam09-orig-sol-01}

\paragraph{Soal.}
Deskripsikan suatu model untuk bidang hiperbolik.

\paragraph{Solusi asli.}
Model setengah bidang atas adalah manifold
\[
 \mathbb H^2=\{(x,y)\in\mathbb R^2:y>0\}
\]
dengan metrik Riemann
\[
 g=\frac{dx^2+dy^2}{y^2}.
\]
Ia lengkap, terhubung, dan mempunyai kelengkungan Gaussian konstan $-1$.
Untuk memeriksa kelengkungan, tulis $g=e^{2u}(dx^2+dy^2)$ dengan
$u=-\log y$. Rumus kelengkungan metrik konformal memberi
\[
 K=-e^{-2u}(u_{xx}+u_{yy})
 =-y^2\left(0+\frac1{y^2}\right)=-1.
\]
Geodesiknya adalah garis-garis vertikal dan busur-busur lingkaran yang
memotong batas ideal $y=0$ secara ortogonal. Transformasi M\"obius nyata
\[
 z\longmapsto\frac{az+b}{cz+d},
 \qquad ad-bc=1,
\]
bertindak sebagai isometri; setelah mengidentifikasi matriks yang berbeda
tanda, grup isometri berorientasi adalah $\mathrm{PSL}_2(\mathbb R)$. Model
cakram Poincar\'e bersifat isometrik dengan model ini melalui transformasi
pecahan linear.

\subsection*{Ujian 10, soal tanpa solusi sumber (O011-EXAM10-ORIG-SOL-01)}
\label{o011-exam10-orig-sol-01}

\paragraph{Soal.}
Jelaskan pasangan istilah ``ekstrinsik dan intrinsik'' dalam konteks manifold.

\paragraph{Solusi asli.}
Suatu konsep bersifat \emph{intrinsik} jika dapat didefinisikan dari manifold
dan struktur yang dibawanya tanpa memilih suatu penanaman ke ruang luar.
Contohnya ialah atlas mulus, ruang singgung, metrik Riemann, koneksi
Levi--Civita, panjang kurva, jarak geodesik, dan tensor kelengkungan Riemann.
Isometri manifold mempertahankan semua data tersebut.

Suatu konsep bersifat \emph{ekstrinsik} jika bergantung pada cara manifold
direalisasikan sebagai submanifold ruang ambien. Untuk
$M\subseteq\mathbb R^N$, medan normal, bentuk fundamental kedua, pemetaan
Weingarten, dan kelengkungan rata-rata memakai ruang ambien dan dapat berubah
ketika penanaman diubah, walaupun metrik intrinsik tetap sama.

Perbedaan ini tidak selalu tampak dari rumus awal. Kelengkungan Gaussian suatu
permukaan pertama-tama dapat dihitung melalui bentuk fundamental kedua, tetapi
Theorema egregium menunjukkan bahwa nilainya sebenarnya ditentukan sepenuhnya
oleh metrik intrinsik. Jadi ``intrinsik'' atau ``ekstrinsik'' merujuk pada
ketergantungan geometris konsep, bukan sekadar pada bentuk rumus yang dipakai
untuk menghitungnya.
